\documentclass[10pt]{article}
% Shared layout for the CV (XeLaTeX). Content lives in cv_en.tex / cv_ru.tex.
\usepackage{fontspec}
\usepackage[a4paper,top=14mm,bottom=14mm,left=16mm,right=16mm]{geometry}
\usepackage{xcolor}
\usepackage{enumitem}
\usepackage[hidelinks]{hyperref}

% NoCommon: Cambria's fi/fl ligatures extract as garbage in the PDF text layer (bad for ATS parsers).
\setmainfont{Cambria}[BoldFont=Cambria Bold,ItalicFont=Cambria Italic,BoldItalicFont=Cambria Bold Italic,Ligatures={TeX,NoCommon}]

\definecolor{linkblue}{HTML}{2F5D8A}
\definecolor{muted}{HTML}{8A8A8A}
\definecolor{todo}{HTML}{C2410C}

\pagestyle{empty}
\setlength{\parindent}{0pt}
\setlength{\parskip}{0pt}
\linespread{1.04}
\fontsize{10pt}{12.6pt}\selectfont

% Placeholder for facts that still have to be filled in. Remove \todo{} before sending.
\newcommand{\todo}[1]{{\color{todo}\textbf{[#1]}}}
\newcommand{\link}[2]{\href{#1}{\color{linkblue}#2}}
\newcommand{\plink}[2]{\href{#1}{\color{linkblue}\textbf{\underline{#2}}}}

\newcommand{\cvsection}[1]{%
  \par\vspace{7pt}{\large\scshape #1}\par\vspace{1pt}%
  {\color{black}\hrule height 0.5pt}\par\vspace{4pt}}

% \entry{Name}{[context]}{Location}
\newcommand{\entry}[3]{%
  \par\vspace{3pt}\textbullet\ \textbf{#1}\ {\color{muted}\itshape #2}\hfill #3\par}
% \role{Role}{Dates}
\newcommand{\role}[2]{%
  \hspace{1.1em}{\itshape #1}\hfill{\itshape #2}\par\vspace{1pt}}

\setlist[itemize]{label={$\circ$},leftmargin=2.4em,itemsep=1.2pt,topsep=1pt,parsep=0pt,partopsep=0pt}
\newenvironment{points}{\begin{itemize}}{\end{itemize}}
\newenvironment{plainlist}{\begin{itemize}[label={\textbullet},leftmargin=1.4em,itemsep=2pt,topsep=1pt]}{\end{itemize}}

\usepackage{polyglossia}
\setdefaultlanguage{russian}
\setotherlanguage{english}
\hypersetup{pdftitle={Сергей Воронов - технический геймдизайнер},pdfauthor={Сергей Воронов}}

\begin{document}

% ---------- Шапка ----------
{\LARGE\textbf{Сергей Воронов}}\hfill \link{mailto:sergei.voronov.workspace.1991@gmail.com}{sergei.voronov.workspace.1991@gmail.com}\par
\vspace{2pt}
{\large\itshape Технический геймдизайнер, Берлин} {\itshape (открыт к удалёнке и релокации)}\hfill
\link{https://sergeivoronov.dev/}{портфолио} $\cdot$
\link{https://www.linkedin.com/in/sergei-voronov-96aa3284}{linkedin} $\cdot$
\link{https://sergeivoronov.itch.io/}{itch.io} $\cdot$
\link{https://store.steampowered.com/app/2058730/The_Darkest_Files/}{steam}\par
\vspace{4pt}
\textbf{Игровые системы} $\cdot$ \textbf{Быстрое прототипирование} $\cdot$ \textbf{Левел-дизайн} $\cdot$ \textbf{Unity} $\cdot$ \textbf{Unreal Engine} $\cdot$ \textbf{Articy: Draft X}\par

% ---------- Профиль ----------
\cvsection{Профиль}
Технический геймдизайнер, \textbf{3 года в геймдеве}: профильное образование, коммерческая стажировка в \textbf{отмеченной наградами студии}
Paintbucket Games и \textbf{13 играбельных проектов} на itch.io. Проектирую и реализую \textbf{игровые системы},
в основном в \textbf{Unity} и \textbf{Unreal}. Мои главные сильные стороны --- \textbf{быстрое прототипирование}
и \textbf{левел-дизайн}. Геймдев --- моя осознанная вторая профессия: после \textbf{9 лет в адвокатуре} я ушёл из
стабильной сферы, чтобы делать игровые системы, которые дарят игрокам незабываемый опыт. Ищу позицию \textbf{Technical Game Designer}, \textbf{Game Systems Designer} или \textbf{Level (Technical) Designer}, где смогу превращать
дизайнерский замысел в работающий геймплей. 

% ---------- Опыт ----------
\cvsection{Опыт}
\entry{Paintbucket Games}{[инди-студия, \plink{https://store.steampowered.com/app/2058730/The_Darkest_Files/}{The Darkest Files: Battalion DLC}]}{Берлин}
\role{Стажёр по техническому дизайну (Unity, Articy: Draft X)}{март 2026 -- июль 2026}
\begin{points}
  \item Сюжет должен был реагировать на находки игрока: я собрал \textbf{логику прогрессии и гейтинга} и конфиги, благодаря чему \textbf{была реализована непрерывная игровая логика}.
  \item Собирал \textbf{сцены, чекпоинты и камеры}, \textbf{тестировал и отлаживал} игровые последовательности ---
        закрыл 6 багов в 3 сценах до дедлайна 31.07.2026.
  \item Был связующим звеном между \textbf{программистами, нарративщиками, техартистами и продюсерами}
        в недельных спринтах (\textbf{Codecks}, \textbf{Plastic SCM}) --- все задачи команды закрывались в текущем спринте.
  \item Рекомендация: \textbf{Йорг Фридрих}, CEO и сооснователь (по запросу).
\end{points}

\entry{Проекты в геймдеве}{[B.A. Game Design и геймджемы]}{Берлин}
\role{Технический / геймплейный дизайнер, соло и в командах из 2--4 человек}{2023 -- н.в.}
\begin{points}
  \item \plink{https://sergeivoronov.itch.io/effacement}{Effacement} \textit{(соло, Unity)}:
        в детективах игрок часто теряется в собственных уликах, поэтому я спроектировал \textbf{систему сбора и связывания
        улик} с \textbf{ветвящимися диалогами} и \textbf{несколькими концовками}, добавил \textbf{повествование
        через окружение} и итерировал по \textbf{плейтестам}, пока тестеры не стали раскрывать дело без подсказок.
  \item \plink{https://prabhdeep-chhatwal.itch.io/berlin-delivery-driver}{Berlin Delivery Driver} \textit{(команда из 3, Unity Visual Scripting)}:
        чтобы трафик и давление окружения ощущались честно, спроектировал \textbf{цикл доставки, оценку маршрутов,
        штрафы и бонусы}, реализовал \textbf{ИИ трафика} и \textbf{настроил сложность}.
  \item \plink{https://shervin-shafizad.itch.io/breath-of-the-dammed}{Breath of the Damned} \textit{(Global Game Jam 2026, команда из 4)}:
        спроектировал \textbf{боевую систему и систему ресурсов} с постоянным дефицитом и провёл уровень
        \textbf{от блокаута до финала} --- играбельный хоррор-прототип за \textbf{48 часов}.
  \item \plink{https://lol-wal.itch.io/what-actually-makes-us-punk}{What ``Actually'' Makes Us Punk?} \textit{(команда из 2)}:
        оригинал был чисто диалоговой игрой --- я добавил \textbf{ритм-систему} (тайминг-окна, проверка ввода, фидбек),
        вплетённую в сюжетные сцены, и дал игре настоящий геймплей, не потеряв её идентичность.
  \item \plink{https://duckmaniac.itch.io/operator}{Operator!} \textit{(команда из 4)}: реализовал единую систему \textbf{«слоты --- провода --- бонусы ---
        входящие звонки»} и \textbf{функциональный интерфейс} для игры об управлении телефонным коммутатором.
  \item \plink{https://sergeivoronov.itch.io/critique-mecontent}{Critique Mécontent} \textit{(соло)}: \textbf{процедурная генерация
        по правилам} абстрактных картин; решения игрока отслеживаются и ведут к \textbf{двум концовкам}.
  \item Также: \link{https://sergeivoronov.itch.io/errate-den-dialekt}{Errate den Dialekt 2.0},
        \link{https://sergeivoronov.itch.io/wasteland-gladiators}{Wasteland Gladiators},
        \link{https://sergeivoronov.itch.io/the-duel}{The Duel},
        \link{https://sergeivoronov.itch.io/room-221}{Room 2.21},
        \link{https://sergeivoronov.itch.io/panzer-pets}{Panzer Pets},
        \link{https://imfawbn.itch.io/cheese-grader}{Cheese-Grater},
        \link{https://ooleks111.itch.io/drunk-driving-simulator-3d}{Drunk Driving Simulator 3D}.
        Регулярно участвую в \textbf{Ludum Dare}, \textbf{Global Game Jam} и \textbf{Berlin Mini Jam}.
\end{points}

\entry{Адвокат по уголовным делам}{}{Воронеж, Санкт-Петербург}
\role{9 лет в защите институтов гражданского общества}{2014 -- 2023}
\begin{points}
  \item Многотомные дела с 5000+ страниц материалов требовали системного подхода: я проводил \textbf{структурный анализ},
        выдвигал и проверял \textbf{гипотезы}, готовил \textbf{документацию} и защищал доверителей в правоохранительных
        органах и суде: \textbf{более 50 уголовных дел прекращено} без судимости, доверители сэкономили \textbf{более 10 млн рублей} на штрафах, которые не пришлось платить.
  \item Каждый день искал \textbf{пограничные случаи и лазейки} в сложных системах правил. Тот же навык я теперь
        применяю, чтобы конструировать и балансировать игровые системы.
\end{points}
\role{Преподаватель уголовного процесса, Санкт-Петербургский институт адвокатуры}{2021 -- 2023}
\begin{points}
  \item Разрабатывал \textbf{учебные материалы} и вёл циклы лекций для аудитории разного уровня подготовки --- \textbf{более 200 слушателей}.
\end{points}

% ---------- Образование ----------
\newpage
\cvsection{Образование}
\entry{B.A. Game Design}{University of Europe for Applied Sciences, Берлин}{2023 -- 2027}
\hspace{1.1em}{\color{muted}\itshape Фокус: прототипирование геймплея, технический дизайн, левел-скриптинг, системы взаимодействия}\par
\entry{Юриспруденция, специалитет}{Воронежский государственный университет}{2008 -- 2014}
\hspace{1.1em}{\color{muted}\itshape Квалификация «юрист»}\par

% ---------- Навыки ----------
\cvsection{Навыки}
\begin{plainlist}
  \item \textbf{Движки}: Unity (C\#, Visual Scripting, 2D / 3D), Unreal Engine (Blueprints, C++)
  \item \textbf{Дизайн}: игровые системы, быстрое прототипирование, левел-дизайн и блокаут, баланс, плейтесты, UI
  \item \textbf{Нарратив и данные}: Articy: Draft X, ветвящиеся диалоги, конфиги и балансные таблицы
  \item \textbf{Контроль версий}: Git (GitHub), Plastic SCM, Fork
  \item \textbf{Продакшн}: Jira, Codecks, Miro; недельные спринты, дейли, внутренние ревью и плейтесты
  \item \textbf{Инструменты}: Realtime-CSG, 3D-моделирование; AI-инструменты (Claude, Codex, Meshy, Tripo)
\end{plainlist}

% ---------- Языки ----------
\cvsection{Языки}
\begin{plainlist}
  \item \textbf{Русский}: родной
  \item \textbf{Английский}: B2 (IELTS 6.5)
  \item \textbf{Немецкий}: A1.2
  \item \textbf{Турецкий}: A1
\end{plainlist}

% ---------- Хобби ----------
\cvsection{Хобби}
\begin{plainlist}
  \item \textbf{Бег}: пробежал \textbf{Берлинский марафон 2025}, медаль финишера.
  \item \textbf{Гиревой спорт}.
\end{plainlist}

\end{document}
